\begin{abstract}
The classification of network traffic records into attack categories is a
core task of intrusion detection. The study evaluated a random forest on a
deduplicated version of the UNSW-NB15 dataset, and compared the forest with
a bagged ensemble of trees, a classification tree and a $k$-nearest
neighbours model. The forest was selected because the forest shares the base
learner of the tree, so that the comparison isolates the effect of
ensembling. Five research questions guided the tuning of the forest. All
pre-processing was fitted within cross-validation folds on a development
partition, and every model was evaluated once on a held-out generalisation
set. The forest attained a macro-averaged $F_1$ measure of 0.607, against
0.591 for the tree, 0.589 for bagging and 0.459 for the nearest neighbour
model, and the improvement over the tree was statistically significant. The
random feature subset produced the entire improvement, because bagging alone
did not outperform the tree. The analysis further established that
$7.1\,\%$ of the records share identical features with a record of a
different category. The conflicts render two attack categories inseparable
and mask most of the advantage of the forest on the separable records.
\end{abstract}
